\scanpage{205}\mspage{150}
(En effet, la formule d'adjonction\deleted{\unsure{ici}}, s'interprète comme un isom. de bifoncteurs, ou ce qui revient au $=$ de foncteurs
\diagram{\begin{tikzcd}[ampersand replacement=\&]
 \unsure{A|\mathcal{A}}'_0 \arrow[r,shift left=2pt,"{\unsure{\varphi|\phi}}"] \arrow[r,shift right=2pt,"\psi"'] \& A_0^{\wedge},
\end{tikzcd}}
or
\[
 \begin{aligned}
 \underline{\mathrm{Hom}}(A'_0,A_0^{\wedge})&\longrightarrow\underline{\mathrm{Hom}}(A',A_0^{\wedge})\\
 \varphi&\longmapsto\varphi\circ\unsure{\rho|p}'
 \end{aligned}
\]
est pleinement fidèle, puisque $\rho'$ est un foncteur de localisation, donc il revient au $=$ de se donner un isom $\varphi\simeq\psi$, ou $\varphi\circ\rho'\simeq\psi\circ\rho'$). On doit donc transformer l'expression
\[
\begin{aligned}
 &\mathrm{Hom}_{A^{\deleted{\unsure{\prime}}}_0}(\underset{\substack{\rule{.35pt}{.4ex}\rule{4em}{.35pt}\rule{.35pt}{.4ex}\\\simeq\,\rho f_!(\xi')}}{f_!^{\unsure{0|\circ}}(\rho'(\xi')},\unsure{\eta|y}_0)
 \underset{\substack{\vert\\\text{comm.}\\\text{de (I)}}}{\simeq}\mathrm{Hom}_{A^{\deleted{\unsure{\prime}}}_0}(\rho f_!(\xi'),\eta_0)\\
 &\underset{\substack{\vert\\\text{adj. }\rho,\unsure{i}}}{\simeq}\mathrm{Hom}_{A_{\deleted{\unsure{0}}}^{\deleted{\unsure{\prime}}}}(f_!(\xi'),\deleted{\unsure{\rho^*}}i(\eta_0))
 \underset{\substack{\vert\\\text{adj. }f_!,f^*}}{\simeq}\mathrm{Hom}_{A'}(\xi',\underset{\substack{\rule{.35pt}{.4ex}\rule{2.8em}{.35pt}\rule{.35pt}{.4ex}\\i'(f_0^*\eta_0)}}{f^*i(\eta_0)})\\
 &\underset{\substack{\vert\\\text{comm. de}\\\text{(II)}}}{\simeq}\mathrm{Hom}_{A'}(\xi',i'(f_0^*(\eta_0)))
 \underset{\substack{\vert\\\text{adj }\rho',i'}}{\simeq}\mathrm{Hom}_{A'_0}(\underset{\substack{\rule{.35pt}{.4ex}\rule{2em}{.35pt}\rule{.35pt}{.4ex}\\\xi'_0}}{\rho'(\xi')},f_0^{\unsure{\alpha|d|*}}(\eta_0)),
\end{aligned}
\]
\deleted{\unsure{ce}} ce qui est la bijection annoncée.

Je résume l'essentiel des résultats obtenus sur $f_!,f_!^{\mathrm{lc}}$ :

\uline{Théorème 2} Soit $f:S'\longrightarrow S$ une flèche \unsure{dans|de} $\mathrm{Cat}$, considérons les foncteurs
\diagram{\begin{tikzcd}[ampersand replacement=\&,column sep=large]
 \deleted{\unsure{f}}\,\mathrm{HOT}_{\unsure{W|w}}(S') \arrow[r,shift left=2pt,"f_!"] \& \mathrm{HOT}_W(S) \arrow[l,shift left=2pt,"f^*"]\\
 \mathrm{HOT}_W^{\mathrm{lc}}(S') \arrow[r,shift left=2pt,"f_!^{\mathrm{lc}}"] \& \mathrm{HOT}_W^{\mathrm{lc}}(S) \arrow[l,shift left=2pt,"f_{\mathrm{lc}}^*"],
\end{tikzcd}}
% archive sheet number: 204
